\chapter*{Обозначения}
\addcontentsline{toc}{chapter}{Обозначения}

{\footnotesize
\setlength{\tabcolsep}{4pt}
\setlength{\extrarowheight}{2pt}
\begin{longtable}{@{}>{\raggedright\arraybackslash}p{0.30\linewidth}>{\raggedright\arraybackslash}p{0.66\linewidth}@{}}
\toprule
Обозначение & Краткое пояснение \\
\midrule
\endfirsthead
\toprule
Обозначение & Краткое пояснение \\
\midrule
\endhead
\midrule
\endfoot
\bottomrule
\endlastfoot
$X$, $Y$, $Z$, \ldots & заглавные латинские буквы обозначают случайные величины. \\
$x$, $y$, $z$, \ldots & строчные латинские буквы обозначают реализации случайных величин. \\
$n$ & объём выборки. \\
$i$ & индекс наблюдения. \\
$j$ & индекс признака или гипотезы. \\
$X^n$ & случайная выборка $(X_1,\ldots,X_n)$. \\
$x^n$ & реализация выборки $(x_1,\ldots,x_n)$. \\
$\prob(\cdot)$ & вероятность события. \\
$\prob(A\mid B)$ & условная вероятность. \\
$F_X(x)$ & функция распределения с.\,в.~$X$. \\
$f_X(x)$ & плотность с.\,в.~$X$. \\
$\mathbb{E}X$ & математическое ожидание с.\,в.~$X$. \\
$\mathbb{D}X$ & дисперсия с.\,в.~$X$. \\
$\mathbb{E}(X\mid Y)$ & условное математическое ожидание. \\
$\mathbb{D}(X\mid Y)$ & условная дисперсия. \\
$\theta$ & неизвестный параметр. \\
$\hat\theta$ & точечная оценка параметра $\theta$. \\
$\bar X_n$ & выборочное среднее. \\
$S_n^2$ & несмещённая выборочная дисперсия. \\
$\mathrm{CI}_{1-\alpha}$ & доверительный интервал уровня $100(1-\alpha)\%$. \\
$H_0$ & нулевая гипотеза. \\
$H_1$ & альтернативная гипотеза. \\
$T(X^n)$ & статистика критерия. \\
$T(x^n)$ & значение статистики на данных. \\
$p$ & p-value. \\
$\alpha$ & уровень значимости. \\
$\beta$ & вероятность ошибки II рода. \\
$1-\beta$ & мощность критерия. \\
$\Phi(x)$ & функция распределения $N(0,1)$. \\
$\phi(x)$ & плотность $N(0,1)$. \\
$N(\mu,\sigma^2)$ & нормальное распределение. \\
$\chi^2_\nu$ & распределение хи-квадрат с $\nu$ степенями свободы. \\
$t_\nu$ & распределение Стьюдента с $\nu$ степенями свободы. \\
$F_{\nu_1,\nu_2}$ & распределение Фишера. \\
\end{longtable}
}
